\section{Introduction}
\label{sec:intro}

Maintenance accounts for a large share of the operating cost of aircraft
engines, and the central decision is deceptively simple: when should an engine
come off the wing? Removing it too late risks an unplanned failure in service,
the most expensive and most dangerous outcome. Removing it too early discards
flight cycles that were already paid for. The classical answer is
\emph{time-based} maintenance: retire every unit at a fixed age chosen from
fleet lifetime statistics~\cite{barlow1960optimum}. It is safe on average but
wasteful, because it treats the longest-lived engine exactly like the
shortest-lived one. \emph{Condition-based} maintenance instead uses each
engine's own sensor data to estimate its remaining useful life (RUL) and act on
it~\cite{jardine2006review,lei2018review}.

The NASA C-MAPSS dataset~\cite{saxena2008damage}, which simulates fleets of
turbofans degrading from a healthy state to failure under noise, multiple
operating regimes and multiple fault modes, has become the standard benchmark
for RUL estimation~\cite{ramasso2014benchmark}. Over the last decade, reported
accuracy on it has improved steadily through convolutional~\cite{babu2016cnn,li2018dcnn},
recurrent~\cite{zheng2017lstm}, hybrid~\cite{aldulaimi2019hdnn} and
attention-based~\cite{ragab2020ats2s} networks. Three gaps nevertheless separate
these models from an autonomous maintenance system.

\emph{First, point estimates are not decisions.} A model that predicts 40
cycles of remaining life with an error of $\pm$5 cycles supports a very
different action from one that predicts 40 cycles with an error of $\pm$40.
Most published C-MAPSS models output a single number. Where uncertainty is
added, typically through Monte-Carlo dropout~\cite{gal2016dropout}, its
calibration is rarely checked, and benchmarks of uncertainty methods for
prognostics show that calibration quality varies widely~\cite{basora2023uqbench}.

\emph{Second, accuracy is measured in the wrong place.} The benchmark protocol
scores each test engine only at the last cycle of a truncated trajectory. This
reveals nothing about how a model behaves as an engine approaches failure,
whether a policy built on it would ground engines in time, or what that policy
would cost relative to simply retiring engines at a fixed age.

\emph{Third, ``digital twin'' is used loosely.} The concept, as articulated for
aerospace by Glaessgen and Stargel~\cite{glaessgen2012digital} and formalized
by Grieves and Vickers~\cite{grieves2017digital}, denotes a virtual model bound
to one specific physical asset and kept synchronized with it. Neural RUL
regressors trained on a whole fleet hold no per-asset state and are not
synchronized with anything, yet they are frequently presented as twins.

At the same time, practical deployment imposes a systems constraint. A model
that runs on or near the engine every cycle must be small and stateless, while
a model that maintains a detailed belief about each engine is naturally hosted
off board, where consulting it costs communication and connectivity. Edge
computing~\cite{shi2016edge} and cascaded inference~\cite{teerapittayanon2016branchynet}
suggest the obvious pattern: run the cheap model always, and escalate to the
expensive one only when needed.

This paper presents \emph{Twin Turbo}, a framework that addresses all three
gaps within that edge/twin pattern. A gradient-boosted edge model runs on every
cycle and is wrapped in Mondrian conformal prediction~\cite{vovk2005alrw,angelopoulos2021gentle},
giving it calibrated intervals at negligible cost. A digital twin, instantiated
once per engine, tracks a sensor-fused health index~\cite{wang2008similarity}
with a particle filter~\cite{arulampalam2002tutorial,orchard2009particle} over
an exponential degradation model whose population prior is learned from
run-to-failure engines~\cite{gebraeel2005residual}. The twin is synchronized
only when the edge model raises one of three triggers (low predicted RUL, a
sharp one-cycle drop, or a periodic resynchronization deadline), and a fixed
policy maps the more cautious of the two calibrated lower bounds to one of four
actions: continue, watch, schedule maintenance, or ground now.

The contributions of this work are:
\begin{enumerate}
  \item \textbf{An instance-specific, synchronizable digital twin} for
  turbofan prognostics, combining a learned health index with Bayesian
  particle filtering of a fleet prior. On the official C-MAPSS test sets it
  attains RMSE between 10.33 and 14.78, and its 90\% posterior intervals are
  empirically calibrated (86--94\% coverage) without tuning.
  \item \textbf{A hybrid edge/twin decision architecture} with conformally
  calibrated edge intervals, trigger-based twin synchronization, and a
  cautious fusion rule that preserves at least 90\% coverage by the union
  bound.
  \item \textbf{A decision-level evaluation protocol} that replays 709
  held-out engines from first cycle to failure and scores every policy by
  failures, unused life and long-run renewal-reward cost rate, against optimal
  age replacement~\cite{barlow1960optimum}, edge-only, twin-only, LSTM and
  oracle baselines, with ablations of every design choice. Under this
  protocol the hybrid policy prevents every failure, reduces cost rate by
  34--47\% relative to age replacement, and consults the twin on 21--28\% of
  cycles.
  \item \textbf{An open, reproducible implementation} with an interactive
  console that replays engines, the fleet and the particle filter live from
  the same traces used for the reported results.
\end{enumerate}

The remainder of the paper is organized as follows.
Section~\ref{sec:related} surveys prognostics methods, uncertainty
quantification, digital twins, edge inference and maintenance decision theory,
and positions this work among them. Subsequent sections describe the dataset
and preprocessing, the edge model, the digital twin and the decision layer,
followed by the experimental protocol, results, ablations and a discussion of
limitations.
